\documentclass[10pt]{article}
\usepackage[a5paper,margin=14mm]{geometry}
\usepackage{fontspec}
\setmainfont{TeX Gyre Termes}
\setsansfont{TeX Gyre Heros}
\setmonofont{TeX Gyre Cursor}[Scale=MatchLowercase]
\usepackage{xcolor}
\usepackage{enumitem}
\usepackage{listings}
\usepackage{fancyvrb}
\lstset{basicstyle=\ttfamily\small,keywordstyle=\bfseries\color{blue},commentstyle=\itshape\color{green!40!black},
  stringstyle=\color{red!70!black},numbers=left,numberstyle=\tiny,numbersep=5pt,frame=single,
  breaklines=true,showstringspaces=false,captionpos=b}
\setlist[itemize]{label=\textbullet,nosep}
\setlist[enumerate,1]{label=(\alph*),leftmargin=*}
\setlist[enumerate,2]{label=\roman*.,leftmargin=*,nosep}

\begin{document}
\section*{Lists and code}
\begin{enumerate}
  \item First choice
    \begin{enumerate}
      \item nested one
      \item nested two
    \end{enumerate}
  \item Second choice\label{item:second}
  \item[\textsc{x}] Custom label
\end{enumerate}
\begin{enumerate}[resume,label=(\alph*)]
  \item Resumed numbering continues here
\end{enumerate}
\begin{description}[leftmargin=2cm,style=sameline,font=\sffamily\bfseries]
  \item[Term] its definition
  \item[Longer term] another definition
\end{description}

\begin{lstlisting}[language=Python,caption={Fibonacci},label=lst:fib]
def fib(n):
    """Return the n-th Fibonacci number."""
    a, b = 0, 1   # start values
    for _ in range(n):
        a, b = b, a + b
    return a
print("fib(10) =", fib(10))
\end{lstlisting}

\lstinputlisting[language=C,firstline=1,lastline=6,caption={From a file}]{snippet.c}

\begin{Verbatim}[numbers=left,frame=lines,commandchars=\\\{\},fontsize=\small,label=Verbatim]
x = 1 + 2 \textbf{\{bold part\}}
y = x ** 2
\end{Verbatim}
Inline: \Verb|\emph{literal}| and \lstinline[language=C]|int main(void)|.
\end{document}
